每小时的请求数、每笔订单的金额、每次接口调用的延迟——三列数摆在面前，要各挑一个分布。翻开手册，二十几个名字排成一张表，每行是密度、期望、方差、矩母函数。表看完了，问题还在原地：这三列数分别该用哪一行？

按公式挑是挑不出来的。密度长什么样是结果，不是理由。把一个分布定下来的是数据被生成的方式：请求由大量互不相干的客户端各自发出，于是计数落进 Poisson；订单金额被少数大客户主导，于是尾巴按幂律拖出去；延迟是排队、重试和超时叠加的产物，既不对称也不轻尾。机制说清楚之后，密度是被推出来的，不必被记住。

所以本章不按字母序列举分布，而按机制分组。同一个机制换一个提问方向就会长出一个新名字：固定试验次数问成功数，得到 Binomial；固定成功数问试验次数，得到 Negative Binomial——两个公式，一套试验。同一个到达过程，数窗口内的事件个数是 Poisson，量一次等待是 Exponential，量第 \(k\) 次等待是 Gamma——三个名字，一件事的三个投影。看清这层结构，要记的东西会少掉一大半。

\begin{center}
\begin{tikzpicture}[x=1cm,y=1cm,font=\footnotesize,>=stealth]
  \foreach \y/\m/\d in {%
    0/{独立重复的成败}/{Bernoulli、Binomial、Geometric、Negative Binomial},
    -1.15/{有限总体、不放回}/{Hypergeometric},
    -2.30/{随机时刻的到达流}/{Poisson、Exponential、Gamma},
    -3.45/{有界比例与概率向量}/{Uniform、Beta、Dirichlet},
    -4.60/{大量微小扰动的叠加}/{Gaussian、\(\chi^2\)、Student \(t\)},
    -5.75/{少数极端观测主导}/{Pareto、Cauchy、Gumbel、Weibull}}{
    \draw[black!45,fill=blue!8] (0,\y-0.4) rectangle (4.3,\y+0.4);
    \node at (2.15,\y) {\m};
    \draw[->,black!55] (4.45,\y) -- (5.15,\y);
    \node[right] at (5.25,\y) {\d};}
\end{tikzpicture}

\emph{左列是机制，右列是它长出来的名字。选分布的动作发生在左列；右列只是记账。同一行里的几个名字往往只差一个提问方向。}
\end{center}

\hyperref[sec:05-Probability/06-Common-Distributions/01]{``成败试验中的 Bernoulli、二项与等待分布''}把一条 0--1 序列摆出来，然后指出四个分布读的是同一条序列。区别全在停止规则：谁是你定死的，谁是随机的。这一节顺带说明为什么在 A/B 测试里``看到显著就停''会让结论失效——停止规则是模型的一部分，改了它，分布就换了。

\hyperref[sec:05-Probability/06-Common-Distributions/02]{``不放回抽样与 Hypergeometric 分布''}拿走独立性。从一个已经标注完的数据集里抽样复核，每抽走一条，剩下的组成就变了。期望不受影响，方差被压缩，压缩的幅度由抽样比例决定。样本相对总体很小时两者几乎重合，这也解释了为什么大多数场合可以放心用二项。

\hyperref[sec:05-Probability/06-Common-Distributions/03]{``Poisson 计数、Exponential 等待与 Gamma 时间''}是本章的主干。一条时间轴上的事件点，数个数、量一次间隔、量到第 \(k\) 个点的总时长，得到三个分布。无记忆性在这里被摆到明面上：它是一个非常强的假设，强到几乎没有真实系统满足。

\hyperref[sec:05-Probability/06-Common-Distributions/04]{``Uniform、Beta 与 Dirichlet 分布''}换一个层次——被建模的对象不再是观测，而是参数本身。Beta 是比例的分布，观测到成功和失败之后它仍然是 Beta，只是参数加了两个计数。这种``更新后不出族''的性质叫共轭，它买来了闭式更新，代价是形状被锁在一个不宽的家族里。Dirichlet 把同一套逻辑搬到多类别，是主题模型和平滑估计的默认先验。

\hyperref[sec:05-Probability/06-Common-Distributions/05]{``Gaussian 与多元 Gaussian''}回答一个很多人跳过的问题：为什么它是默认？两个彼此独立的理由——大量微小扰动的和会趋向它，以及在只知道均值和方差的前提下它是最不承诺的那个分布。多元情形里协方差矩阵完全决定了等高线椭球的朝向和粗细，这一步直接接上谱定理。

\hyperref[sec:05-Probability/06-Common-Distributions/06]{``从 Gaussian 到卡方与 Student \(t\)''}讲的是构造出来的分布。卡方是标准正态的平方和，\(t\) 是正态除以一个独立的卡方尺度。它们之所以在统计推断里到处出现，是因为你用同一批数据同时估计了均值和方差——分母也在抖，尾巴因此比正态厚。

\hyperref[sec:05-Probability/06-Common-Distributions/07]{``重尾分布：当均值与方差失效时''}是本章判断价值最高的一篇。均值和方差是积分，积分可以发散：Cauchy 的均值不存在，Pareto 的哪些矩存在取决于尾指数。这不是数学奇观——大数定律和中心极限定理的前提在这里塌掉，样本均值不收敛，基于它的置信区间失去意义。收入、文件大小、词频、请求延迟都活在这一档里，这也是``平均延迟''作为指标为什么长期误导人。

\hyperref[sec:05-Probability/06-Common-Distributions/08]{``极值分布与尾部风险''}把问题从``平均而言''换成``最坏的一次''。块最大值经过规范化之后只有三种可能的极限分布，具体落进哪一种由母体的尾部性状决定，建模者说了不算。关心峰值时用正态近似会系统性地低估风险，这一节给出替代路线和它的失效边界。

首次通读建议按顺序走前六篇，重点放在机制配对上：固定次数与固定成功数、放回与不放回、计数与等待、观测与参数、单变量与协方差几何。最后两篇不是附录。只要你的风险由少数极端观测决定——容量规划、成本控制、SLA 承诺、金融敞口——就应该优先读它们，并且在计算任何均值或方差之前先确认这两个量存在。

读完之后可检验的变化，是拿到一列数时的第一反应从``用哪个公式''变成三个问题：这些数是怎么产生的？我固定了什么、什么在随机波动？极端值出现的频率高到会让平均失效吗？答不上来的时候，公式记得再熟也没有落脚点。
